\begin{table}[H]
\centering
\footnotesize
\renewcommand{\arraystretch}{1.4}
\rowcolors{2}{gray!10}{white}
\begin{tabularx}{\textwidth}{>{\raggedright\arraybackslash}p{0.20\textwidth} >{\raggedright\arraybackslash}p{0.225\textwidth} >{\raggedright\arraybackslash}X}
\toprule
\textit{Command} & \textit{Arguments} & \textit{What it does} \\
\midrule
\textbf{WHERE} & boolean predicate & Filters individual rows before grouping. Cannot see aggregates, window functions, or select aliases. \\
\textbf{GROUP BY} & column or expression list & Collapses rows to one output row per distinct key combination. \\
\textbf{HAVING} & boolean predicate & Filters whole groups after aggregation. The only legal home for an aggregate predicate. \\
\textbf{DISTINCT} & applies to the select list & Removes duplicate rows compared across the entire select list. Treats two \texttt{NULL}s as equal. \\
\textbf{DISTINCT ON} & \texttt{(expr, ...)} & Postgres. Keeps one row per distinct value of the expressions, chosen by the leading \texttt{ORDER BY} terms, which must match. \\
\textbf{ORDER BY} & \texttt{expr [ASC|DESC] [NULLS FIRST|LAST]} & Sorts the output. Postgres places \texttt{NULL}s last under \texttt{ASC} and first under \texttt{DESC}. \\
\textbf{LIMIT / OFFSET} & \texttt{n} & Caps the number of rows returned, or skips the first \texttt{n}. \\
\textbf{WITH} & \texttt{name AS (query)} & Common table expression. Names a subquery for reuse. \texttt{WITH RECURSIVE} walks hierarchies. \\
\textbf{CASE} & \texttt{WHEN cond THEN val [ELSE val] END} & Row-level branching. Returns the first match, or \texttt{NULL} when nothing matches and \texttt{ELSE} is absent. \\
\textbf{FILTER} & \texttt{(WHERE cond)} after an aggregate & Restricts which rows the aggregate consumes. Postgres shorthand for \texttt{SUM(CASE ...)}. \\
\textbf{EXISTS / NOT EXISTS} & subquery & Semi-join and anti-join. Null-safe, unlike \texttt{IN} and \texttt{NOT IN}. \\
\bottomrule
\end{tabularx}
\caption*{Query clauses and structure}
\end{table}

\vspace{-2ex}
